\section{Introduction}

\paragraph{Classical information gain and active learning.}
In many active learning and Bayesian experimental design problems, the goal is not simply to collect observations, but to reduce uncertainty in a way that is useful for a downstream task. Classical information-gain objectives provide a useful measure of uncertainty reduction and have played a central role in Gaussian process active learning, Bayesian optimisation and experimental design. However, these objectives typically treat all parts of the domain uniformly. In many practical settings, uncertainty in some regions may be substantially more important than uncertainty elsewhere, suggesting that a more task-aware notion of information gain may be desirable.

\paragraph{Transductive Active Learning and task-aware exploration.}
Recent work on Transductive Active Learning (TAL) replaces the global objective with a target-dependent quantity
\begin{equation*}
\max_{X \subseteq S,\ |X|\leq T}
I(f_A;y_X),
\end{equation*}
which replaces the global objective with a target-dependent notion of information gain. Rather than measuring information acquired about the entire domain, TAL measures only the information gained about a specified target set. This allows exploration to focus on regions that are relevant to the downstream task and can lead to substantially sharper uncertainty bounds.

Despite this improvement, TAL still imposes a binary notion of relevance: points either belong to the target set or they do not. All locations within the target set are treated equally.

\paragraph{Motivating example.}

Consider an autonomous racing car learning the dynamics of a racetrack. Classical information-gain objectives encourage the vehicle to reduce uncertainty across the entire track, since all regions contribute equally to the objective. If the goal is simply to learn the track as accurately as possible everywhere, this behaviour is reasonable.

Suppose, however, that only a particular sequence of corners is safety-critical or performance-critical. In this setting, it is not obvious that the vehicle should continue allocating samples uniformly around the entire track. One may instead wish to prioritise uncertainty reduction in the important section while still accounting for information gained elsewhere.

TAL partially addresses this issue by allowing one to specify the critical section as a target set. However, all locations inside that section remain equally important. In practice, one might care most about a particular section of the corner, such as the point at which the vehicle begins braking, the turn-in point, or the point of maximum curvature, suggesting a more graded notion of task relevance.

\paragraph{Research question.}

This motivates the following question: can information gain be adapted to account for varying levels of importance across the domain while retaining useful theoretical structure?

Rather than specifying a binary target set, we seek a weighted notion of information gain in which different regions contribute according to a user-defined importance function. In this sense, the weighted setting may be viewed as a continuous analogue of the TAL framework.

\paragraph{A difficulty: loss of submodularity.}

A natural approach is to introduce weights directly into the marginal information gains, giving more importance to uncertainty reduction in some regions of the domain than in others. However, this modification fundamentally changes the structure of the objective.

For ordinary information gain, the objective can be written as a submodular set function, and the corresponding marginal gains satisfy a diminishing returns property. This structure underlies many of the theoretical guarantees in active learning and Bayesian experimental design, in particular allowing us to compare greedy selection objectives directly with maximum informatrion gain.

In the weighted setting, the contribution of an observation depends both on its weight and on its position in the observation sequence through the posterior variance. Consequently, the objective becomes order-dependent and can no longer be represented as a submodular set function. This prevents the direct application of the standard information-gain and transductive active learning framework, motivating the search for an alternative analytical approach.

\paragraph{Main contributions.}

The loss of submodularity means that weighted information gain cannot be analysed using the standard machinery underlying information gain and transductive active learning. Our main observation is that, despite this difficulty, piecewise constant weights admit a natural decomposition into regional information-gain terms.

\begin{proposition*}[Informal split bound]
Suppose the weight function is piecewise constant,
$$
w(x)
=
\sum_{i=1}^m
c_i \mathbf 1_{A_i}(x),
$$
for a partition
$
\mathcal X
=
\bigcup_{i=1}^m A_i.
$
Then the weighted information gain can be bounded by a sum of ordinary information-gain terms on the individual regions, each evaluated with an effective noise level proportional to $\sigma^2/c_i^2$.
\end{proposition*}

This result reduces the analysis of weighted information gain to the analysis of standard information gain on individual regions. Consequently, existing information-gain bounds for Gaussian process kernels can be applied separately on each region, allowing the effects of geometry, volume and intrinsic dimension to be captured explicitly. In particular, the resulting bounds reveal situations in which weighted exploration can be substantially more efficient than global exploration, especially when highly weighted regions occupy only a small subset of the domain or lie on lower-dimensional structures.

\begin{proposition*}[Informal consequence for Squared Exponential kernels]
For Squared Exponential kernels, applying the split bound reveals that restricting attention to smaller regions primarily affects lower-order geometric constants. In particular, regional volume enters only logarithmically into the resulting information gain bounds. Consequently, the largest improvements arise when highly weighted regions have lower effective dimension than the ambient domain.
\end{proposition*}

\begin{proposition*}[Informal consequence for Mat\'ern kernels]
For Mat\'ern kernels, applying the split bound reveals a substantially stronger dependence on regional geometry. Both the volume of a region and its effective dimension influence the leading-order information gain rate. Consequently, concentrating weight on small or lower-dimensional regions can improve not only constants but also the polynomial scaling of the information gain bound.
\end{proposition*}

The remainder of the paper formalises this decomposition, establishes weighted information-gain bounds for piecewise constant weight functions and investigates how regional geometry and intrinsic dimension affect the resulting rates for linear, Squared exponential and Matérn kernels.

